\begin{env}[7.3.13]
\label{IV.2.s7.3.13-fr}
Dado un anillo semilocal noetheriano $A$, llamaremos \emph{fibras
formales} de $A$ a las fibras del morfismo canónico
$f:\operatorname{Spec}(\widehat A)\to\operatorname{Spec}(A)$; para todo
ideal primo $\mathfrak p$ de $A$, la fibra formal en $\mathfrak p$ es,
por tanto, el esquema
$\operatorname{Spec}(\widehat A\otimes_A \boldsymbol{k}(\mathfrak p))$; como el
completado del anillo semilocal $A/\mathfrak p$ es
$\widehat A/\mathfrak p\widehat A=(A/\mathfrak p)\otimes_A\widehat A$,
se ve que la fibra formal de $A$ en $\mathfrak p$ es también \emph{la
fibra formal de $A/\mathfrak p$ en el punto genérico \textup{(0)} de
$\operatorname{Spec}(A/\mathfrak p)$}.
\end{env}

\oldpage[IV-2]{196}

Definida la propiedad $\boldsymbol P$ como en (7.3.1), diremos que
\emph{las fibras formales de $A$ tienen la propiedad $\boldsymbol P$},
o que $A$ es un \emph{$\boldsymbol P$-anillo}, si, para todo
$x\in\operatorname{Spec}(A)$, se cumple
$\boldsymbol P(f^{-1}(x),\boldsymbol{k}(x))$. Puesto que $f$ es \emph{plano}, esto
equivale a decir que $f$ es un $\boldsymbol P$-\emph{morfismo}.
